\textbf{Explanation of samples}


In the first sample the sweater is 13 loops in circumference. Thus the centered
pattern is placed at loops 6, 7 and 8. There's space for another pattern in either
direction at loops 3, 4, 5 and 9, 10, 11. There's not enough space to place two more,
and a single pattern would make things asymmetric. Thus loops 1, 2, 12 and 13 are empty,
so the answer is 4.

\begin{figure}[!ht]
    \centering
    \includegraphics[width=0.15\textwidth]{knittingpatterns_1_pattern}
    \caption{A possible pattern for sample 1}
\end{figure}

\begin{figure}[!ht]
    \centering
    \includegraphics[width=0.65\textwidth]{knittingpatterns_1_full}
    \caption{Applying the pattern to the full width for sample 1}
\end{figure}

In the second sample the sweater is 16 loops in circumference. The first pattern is placed
at loops 7, 8, 9 and 10. Two more are placed at 3, 4, 5, 6 and 11, 12, 13, 14. This leaves
1, 2, 15 and 16, which exactly fits one more pattern that will be perfectly centered at the
back of the sweater, creating no asymmetry. Thus that pattern is placed, leaving no empty
space.

\begin{figure}[!ht]
    \centering
    \includegraphics[width=0.2\textwidth]{knittingpatterns_2_pattern}
    \caption{A possible pattern for sample 2}
\end{figure}

\begin{figure}[!ht]
    \centering
    \includegraphics[width=0.8\textwidth]{knittingpatterns_2_full}
    \caption{Applying the pattern to the full width for sample 2}
\end{figure}